\chapter[Les attentes internes]{Les attentes internes: inscrire Cujas dans une stratégie nationale}

\lettrine{L}'ère des bibliothèques numériques isolées, conçues comme de simples \enquote{répertoire de PDF} ou de données isolées les uns des autres est maintenant révolue. L'essor de la Science Ouverte et la nécessité de valoriser la recherche publique, imposent désormais de nouvelles exigences de découvrabilité ainsi qu'une refonte des infrastructuress documentaires. Le constat technique effecuté sur CujasNum, marqué par l'obsolescence de son logiciel XTF, l'impossibilité d'un moissonnage OAI-PMH par Gallica et la vulnérabilité d'un système reposant sur une compétence humaine unique, met en exergue son dysfonctionnement matériel. Loin d'être isolé, il témoigne des limites structurelles du deuxième âge du numérique, caractérisé par une interopérabilité \enquote{par copie et conversion} de notices plates ou de pertes de données en passant du format MARC au DUblin Core, est devenu inapte à exprimer la richesse sémantique des collections des institutions patrimoniales. Face à ce constat la BIU Cujas a choisi d'inscrire la refonte de sa bibliothèque numérique dans une stratégie nationale. Cette refonte a pour objectif de s'inspirer du Web de données afin de s'appuyer sur les réseaux de bibliothèquesn umériques et de valoriser ainsi ses nombreuses collections.

\section{Du silo local au catalogage partagé}
La refonte de la bibliothèque numérique amène actuellement la BIU Cujas a réfléchir à la manière dont celle-ci travaille. Cette refonte entraîne également une refonte de la chaîne de numérisation, fonctionnant principalement de façon isolée et manuelle. Le DPCN cherche un moyen de sortir de ce carcan afin de pouvoir automatiser d'avantage de tâches en minimisant l'impact humain. Bien sûr, il est important de ne pas supprimer complètement la supervision humaine, mais cela ne les empêche pas de réfléchir à un moyen d'enrayer ce travail devenu avec le temps chronophage.

\subsection{Les limites du catalogage en silo et la dérive des notices déconnectées}
\subsubsection*{Le constat de l'existant}
Actuellement, la création des métadonnées sous CujasNum repose sur une production segmentée des tâches: le prestataire saisie les métadonnées dans la Dublin Core en s'appuyant sur le bordereau fourni au moment de la livraison par la bibliothèque. Une fois, la numérisation effectuée, le technicien en charge du contrôle qualité ouvre la Dublin Core et effectue des micro-corrections lorsque cela est nécessaire. La personne fait attention à ce que le schéma soit bien déclairé, que ce soit les bons guillemets. 

Historiquement, avant 2013, la bibliothèque Cujas ne disposait d'aucun agent formé au catalogage du livre ancien. Il n'existait pas à cette époque de standardisation afin d'indiquer comment la notice devait être construite. La logique était d'utiliser lorsqu'elle existait la notice cataloguée en ligne. Quand il n'y avait pas de modèle existant, la personne chargée de ce catalogage descendait dans le magasin pour décrire la notice à partir de l'ouvrage papier en renseignant autant d'information qu'il lui était possible de le faire. Cette manière de procéder a crée des notices numériques inégales. Certaines étaient riches et d'autres extrêmmeent pauvres ou incomplètes demandant aujourd'hui unne reprise complète de ces notices avant de pouvoir envisager le moindre enrichissement de notices.

En effet, pour éviter que de créer des notices bibliographiques déjà existantes, il est nécessaire pour les bibliothèques de sortir de leur production en silo correspondant au \enquote{stockage de l'information dans une base locale, isolée, dont le catalogue traditionnel est l'archétype}. Le Web de données a pour but de \enquote{sortir l'information de ces silos} pour la rendre exploitable de manière distribuée. L'enrichissiment des notices est possible grâce au web sémantique qui \enquote{permet d'exposer des données, des contenus d'une manière beaucoup plus riche qu'à l'heure actuelle, et permet d'encourager des collaborations transdisciplinaires et le croisement entre professions de la culture ou de la recherche}. 

Pour permettre ces enrichissements, de nombreuses institutions patrimoniales ont effectué des chantieres de nettoyage de notices afin de permettre des alignements sur des référentiels d'autorité comme IdRef. C'est une étape jugée primordiale par celle-ci afin d'avoir des notices propres et enrichies participant à la visibilité de leur bibliothèque numérique par les autres institutions permettant ainsi à celle-ci de pouvoir créer des rebonds en s'appuyant sur ce référentiel. 

En 2023, la bibliothèque Cujas faisait face au problème de déconnexion entre le catalogage et le traitement numérique. Lors de la création des notices bibliographiques, deux méthodes se chevauchent soit les agents créent la notice de toute pièce soit ils se contentent d'un lien cliquable. Il est nécessaire de résoudre ce problème qui impacte sérieusement le reste de la chaîne de numérisation. En effet, \enquote{le catalogage des documents numérisés doit être effectué en début de chaîne : tout retard impacte l'ensebmle des opérations qui suivent}. Un autre problème vient de la politique de catalogage elle-même qui comme nous l'avons vu n'est pas \enquote{bien définie : parfois une notice de document électronique a été créée, parfois un lien vers CujasNum a été ajouté dans le champs Unimarc 856}. Le moissonnage de ces notices étant imparfaits car la notice créée pour CujasNum est celle se retrouvant dans le catalogue de la bibliothèque. Face à ces problèmes répétés, la bibliothèque a pris énormément de retard dans son catalogage et son indexation. En effet, on peut le voir avec les cours polycopiés \enquote{numérisés en 2020 qui n'ont pas été catalogués.}. Un travail rétrospectif est nécessaire une fois que les notices ont été modifiés, de \enquote{modifier rétrospectivement les métadonnées en Dublin Core}. Ces inconvénients créent des retards voir des doublons dans l'outil de découverte Primo. 

La majorité des bibliothèques utilise le Dublin Core simple, constitué de 15 balises mais la plupart des bibliothèques n'en utilise réellement que cinq (creator, identifeir, title, date, type). En effet, cela comprend \enquote{74 \% des fournisseurs qui ne dépassaient pas l'utilisation de cinq éléments}. L'adoption massive de ce format par les institutions patrimoniales comme les bibliothèques a entraîné un appauvrissement des données contrairement aux données que celle-ci pouvait encoder au format MARC. L'appauvrissement constaté des notices bibliographiques résulte paradoxalement de la volonté des bibliothèques de rendre leurs données interopérables et visbiles sur le Web. Face à la diversité et à la complexité des formats d'encodage (tels que le format MARC utilisé dans des systèmes comme Alma), l'interopérabilité directe s'avère difficile à mettre en oeuvre. Les institutions patrimoniales ont donc recours au Dublin Core simple, un standard minimaliste adopté comme \enquote{le plus petit commun dénominateur}. Couplé au protocole OAI-PMH pour le moissonnage automatisé, ce format pivot permet une diffusion immédiate, universelle et à bas coût des données sur le Web, au prix d'une simplification transitoire de l'information bibliographique. En effet, les universités possédant peut de moyens doivent trouver des alternatives à bas coût permettant leur exposition sur le web de données. C'est ainsi que l'idée derrière l'utilisation de ce format est \enquote{de favoriser la création d'archives ouvertes institutionnelles, et de créer un réseau pour les fédérer, et en même temps participer à améliorer la qualité des métadonnées échangées}. Le protocole OAI-PMH assure cet \enquote{échange de métadonnées, à un moment où l'inflation de contenus sur le Web} explose. C'est également une plus-value pour les chercheurs car cela permet de présenter les informations bibliographiques de la même manière ne perturbant nullement l'utilisateur lorsqu'ils consultent plusieurs bibliothèques numériques différentes. 

Toutefois, cette simplification transitoire s'est fait au détriment de la richesse des données. Comme le rappelle Olesea Dubois, le signalement devrait être pensé comme un véritable \enquote{passport du document} exigeant une description exhaustive (aspects physiques, particularités d'exemplaires) essentielle pour les chercheurs. De surcroît, cette logique de diffusion est paralysée pour la bibliothèque Cujas : l'entrepôt OAi-PMH de CujasNum, basé sur une version d'XTF non mise à jour, est devenu impossible à moissonner par les partenataires nationaux, comme la BnF, l'enfermant dans un véritable silo technique. 

Les équipes de Cujas espèrent qu'avec la refonte de leur bibliothèque numérique, les problèmes de moissonnage et de constitutions de notices soient rectifiées afin d'avoir une meilleure visibilité sur le web. 
\subsubsection*{Les exigences internes pour la refonte}
Il est primordial pour les agents de Cujas de parvenir à intégrer leurs notices bibliographiques dans le catalogue SUDOC ce qui n'est pas le cas pour le moment. En intégrant ce catalogue, cela leur permet ainsi de mieux signaler leurs collections de la même manière que le font les autres bibliothèques. 

La majorité des bibliothèques se sont mis à \enquote{rationnaliser la production des données} afin de profiter de ce qu'offre le Web de doonées aux bibliothèques. Il offre l'opportunité de répartir l'activté de description de leurs fonds de manière décentralisée, \enquote{différents acteurs enrichissant le même graphe}. Cela justifie la proposition de la bibliothèque Cujas de \enquote{traiter le catalogage et l'indexation sémantique en amont dans  le SUDOC} pour que la notice \enquote{bavarde} redescende proprement dans Alma et Codex, plutôt que de devoir la corriger à la main au moment du contrôle qualité. Cette proposition formulée par le Département de l'informatique documentaire (DID) et le DPCN est \enquote{d'abandonner la double saisie locale}, entraînant sans arrêt le retour en arrière sur une tâche en ayant l'impression qu'elle n'aboutit jamais ou plus tardivement que le souhait initial. Il faudrait que dans le futur le prototype Codex puisse s'interconnecter nativement avec le SUDOC pour récupérer des notices propres et éviter les doublons. 

Une autre des volontés des experts métiers est d'avoir accès à un back-office structuré où à l'intérieur sont prédéfinis des rôles pour chaque profil (catalogueur, éditeur, responsable). Cela permettrait d'avoir une infrastructure balise où les droits d'accès sont délimités à l'avance en fonction des tâches assignées à cet agent sans mettre en péril le reste de la production dans le cas où une nouvelle connexion par un autre agent ne destabilise la structure. Les profils permettrait ainsi à la responsable de la bibliothèque numérique de pouvoir assigner des tâches comme lorsqu'ils souhaiteront faire de l'éditorialisation ayant accès uniquement au partie dédiée permettant de le faire. Cela aurait comme conséquence de décentraliser le travail et de sortir de la dépendance sur les compétences du responsable informatique. La bibliothèque aimerait que le back-office soit \enquote{administrable par un bibliothécaire}, suffisant balisée et facile à prendre en main afin de décharger le responsable DIT qui s'occuperait de la partie maintien technique et mises à jour de la bibliothèque, là où les experts métiers pourraient s'occuper totalement de la partie production, mise en ligne et correction des micro-erreurs typographiques. 

Ce n'est pas la seule exigence interne des agents en charge de la bibliothèque car avant de mettre en place ces droits, il est tout aussi important de s'occuper de leurs métadonnées. Cela représente un chantier important à réaliser en priorité sinon la structure de la bibliothèque n'aurait pas été nettoyé et ne pourra donc pas être enrichi. Cujas est loin d'être la seule bibliohèque patrimoniale à avoir rencontré ce problème lié aux métadonnées. En effet, toutes les institutions patrimoniales ayant souhaité une refonte de leur bibliothèque numérique sont passées par cette étape jugée comme importante par des experts comme Pauline Rivière ou Olesea Dubois. Au cours de l'entretien de Pauline Rivière, elle a insisté en décrivant cette \enquote{étape comme étant indispensable}. Elle en a eu recours afin de \enquote{lier ces 5 000 notices de la Bibliothèque Sainte-Geneviève (BSG)}. Pour réaliser cet alignement, elle a utilisé OpenRefine qui lui a permis d'aligner ces notices bibliographiques au référentiel d'IdRef. Olesea Dubois mentionne également l'importance de nettoyer ces notices car \enquote{si on ne les nettoie pas, cela crée énormément de bruit, créant ainsi des résultats insuffisants}, alors que lorsqu'elles \enquote{sont normés et standardisée, cela ne représente pas un problème}. 

Cette reprise est nécessaire, car sinon cela crée le problème que l'on observe actuellement sur CujasNum. En se rapprochant del a standardisation, c'est mieux et \enquote{cela permet d'éviter de nombreux problèmes}. Pauline Rivière explique qu'il faut faire attention pour les notices récentes, \enquote{les matchs étaient bons, mais ils ont dû repasser à la main sur les autres pour les vieux ouvrages}. En effet, comme expliquée en amont, les notices bibliographiques de livres anciens pour la plupart peuvent être pauvres car il n'existait pas avant les années 2010 de format de standardisation propre pour ce type d'ouvrage. Toutefois, elle ne regrette pas d'avoir fait ce choix.

Les alignements permettent d'intégrer les notices bibliographiques dans le Web de données amenant une meilleure visibilité de celles-ci sur Internet et permettant aux bibliothèques de lutter contre \enquote{l'info-obésité} à laquelle sont confrontés quotidiennement les usagers. 
\section{Le Web de données et les nouveaux protocoles}
Avant la diffusion des données était davantage passive où les bibliothèques numériques interrogeait les entrepôts OAI afin de moisssonner leurs idées, sans forcément intéragir directement avec celle-ci. La transition de la bibliothèque numérique de la BIU Cujas vers le Web sémantique ne saurait être réduite à une simple mise à jour d'infrastructure : elle engage une véritable refondation de l'objet documentaire lui-même. Aujourd'hui, le document numérique rompt définitivement avec la logique de faire la numérisation sous forme de captures d'image. Désormais, la représentation du document numérisé est devenu plus complexe. Comme l'a théorisé Jean-Michel Salaün au sein du collectif Roger T. Pédauque, \enquote{il est probable que le document numérique donne une dimension nouvelle à la notion d’intertextualité. Celle-ci se prolonge dans l’inter-sémioticité. L’inter-sémioticité du document numérique qui intègre texte, écrit, image et son ne change pas la nature du texte, elle en accroît en revanche la complexité. Les rapports d’inclusion réciproque constituent un objet de recherche nouveau qui devrait mieux spécifier la relation entre sémantique et sémiotique}. Le document numérique est devenu un objet structurée en trois dimension : à travers la forme, le texte et l'interconnexion sémantique. C'est pourquoi, faire entrer le patrimoine juridique dans le Web de données (Linked Open Data) consiste à matérialiser cette déconstruction pour substituer au document-objet traditionnel un objet dynamique et ouvert d'entités liées.

\subsubsection{La rigueur de la structure logique}
La structure logique des documents numérisés constitue le fondement de toute bibliothèque numérique patrimoniale. Comme le rappelle Isabelle Westeel, à partir des années 1970, l'émergence historique des langages de balises vient \enquote{de la nécessité de structurer l'information}, en se préoccupant \enquote{du contenu du document et non pas de sa mise en forme (contrairement au HTML qui gère des balises de mise en forme)}. Cette évolution consacre une distinction majeure dans la composition du document entre \enquote{une structure logique (contenu organisé) et d'une structure physique (mise en forme de l'information)}. Toutefois, la simple conformité syntaxique n'est pas suffisante pour garantir l'interopérabilité des données. Le document est \enquote{bien formé}, lorsqu'il respecte \enquote{l'ensemble des règles de la recommandation} ou l'utilisation de la DTD \textit{Document Type Definition} et des schémas XML. La DTD liste les éléments, les attributs associés dans le document, participant l'interopérabilité des métadonnées.

Ce sont ces mêmes métadonnées qui sont moissonnées à travers les entrepôts d'OAI-PMH des bibliothèques. Actuellement, l'entrepôt d'OAI-PMH de CujasNum rencontre des problèmes de moissonnages. L'une des raisons vient de l'absence de balises manquantes obligatoires comme \enquote{earliestDatestamp}, bloquant le moissonnage automatique par Gallica. Face à ce problème récurrant auquel est confronté la bibliothèque Cujas, l'établisssement souhaite qu'avec la nouvelle bibliothèque le moissonnage se fera de nouveau. En effet, à cause de l'impossibilité de faire moissonner leurs collections par les institutions nationales, comme la BNF, cela entraîne un manque de visibilité de leurs documents dans le Web de données.

\subsection{Le dépassement de l'OAI-PMH : de la diffusion passive aux données liées}
Depuis les années 2000, le protocole OAI-PMH (\textit{Open Archives Initiative Protocol for Metadata Harvesting})couplé au format Dublin Core simple a constitué la brique standard incontournable pour le moissonnage des métadonnées sur le Web. Si cette simplification transitoire a permis une diffusion à bas coût de l'information bibliographique, elle se heurte aujourd'hui aux exigences scientifiques des bibliothèques de recherche patrimoniale. Comme le démontre Vincent de Lavenne de la Montoise, le Dublin Core simple, limité à quinze éléments, est incapable de reproduire la richesse des documents complexes et entraîne \textit{de facto} une disparition des liens d'autorité. En effet, le Dublin Core simple ne possède \enquote{aucun mécanisme permettant de représenter ces notices d'autorité qui concerne des entités non documentaires (personnes, concepts, etc.)}. Les notices d'autorité sont la base permettant de créer des liens entre les notices bibliographiques. A la bibliothèque Cujas, cette \enquote{mise à plat} imposée par ce format provoque une destruction de l'arborescence initiale des documents complexes. Un périodique de doctrine, par exemple, s'aritcule obligatoirement selon une arborescence stricte  (titre de la revue, année de publication, fascicule, article) que le système actuel peine à restituer de manière interopérable, empêchant ainsi la bibliothèque de mettre en avant ce fonds numérisés.

Vincent de Lavelle rappelle que ce format a été adopté massivement, environ \enquote{74 \% des forunisseurs ne dépassaient pas l'utilisation de cinq éléments sur les quinze qui sont dc: creator, dc: identifier, dc: title, dc: date et dc: type}, participant à \enquote{l'exclusion d'enrichissement}. La souplesse qu'offre le Dublin Core inquiète les fournisseurs de données face à cette perte de l'information, amenant à \enquote{concacténer des données diverses dans des champs textuels}. La balise <dc: publisher> en est le parfait exemple, car il est souvent arriver au bibliothèque d'y concaténer le nom de l'éditeur et le lieu de la publication. Il faut donc envisager d'autres moyens de décrire les documents sur le Web en ayant, par exemple, recours au format de sérialisation.

\subsubsection{Décrire les documents sur le Web}
Si le protocole OAI-PMH a permis de diffuser massivement des métadonnées sous forme de notices plates, la transition vers le Web sémantique impose de repenser la description même des ressources. Décrire les documents sur le Web ne consiste plus à rédiger une notice bibliographique isolée, mais à insérer chaque entité (personne, concept) dans un réseau mondial de données sémantiques. 

Cette transformation s'appuie sur le modèle RDF (\textit{Resource Description Framework}), qui s'affranchit de la rigidité des bases de données traditionnelles (SQL ou XML) pour modéliser l'information d'un graphe décentralisé composé de triplets (Sujet, Prédicat, Objet). Dans cette architecture sémantique, la description, par exemple, d'un traité de la listet Pfister-Roumy devient malléable et évolutive: l'ajout de nouvelles propriétés ou de liens vers des référentiels externes s'effectue par simple incrémentation de triplets, sans risque de briser la structure logique du système d'information. Cela permettrait de mieux s'adapter aux pratiques des chercheurs, en termes de données. En 2017, les spécialistes d'ingénierie documentaire s'interrogeait sur le format de données à proposer aux chercheurs ainsi que sur leurs pratiques. Par ailleurs, la bibliothèque Cujas s'interroge également dessus souhaitant davantage proposer un format de données correspondant aux pratiques de recherche des chercheurs. Antoine Courtin plaidait pour un grand réalisme pragmatique en rappellant que \enquote{les pratiques des chercheurs, en termes de données, se limitent bien souvent à copier des données dans un document de traitement de texte. Dans ce contexte, est-il opportun de chercher à adopter des formats, des langages de requête, en somme des logiques qui ne sont pas ceux des communautés que l'on sert?}. Finalement, au lieu d'imposer des formats sémantiques trop complexes et rigides aux chercheurs, les bibliothèques doivent s'adpater à leurs gestes quotidiens les plus simples, comme le copier-coller de texte ou de référence. Ils cherchent avant tout à manipuler des données brutes. 

L'usage des technologies du Web sémantique a un potentiel fort d'amélioration des services rendus aux usagers. Ces technologies ouvrent la voie à l'adoption d'autres standards notamment le protocole IIIF permettant la multi-diffusion des images. 

\subsection{Le protocole IIIF et la multi-diffusion}
Grâce au protocole IIIF (\textit{International Image Interoperability Framework}), la bibliothèque Cujas pourrait reconstruirre virtuellement ses collections, comme la liste Pfister-Roumy, en appellant les manifestes d'images stockés à la BNF, offrant une alternative scientifique robuste aux \enquote{Essentiels du Droit} de Gallica sans héberger physiquement les fichiers. En effet, c'est l'une des ambitions de l'établissement, à travers cette refonte, elle envisage de devenir une alternative aux documents de droit exposés sur Gallica. L'attente interne n'est plus de se contenter d'exposer passivement ses notices bibliographiques sur sa bibliothhèque numérique, mais d'interconnecter activement ses notices dans le système documentaire. En adoptant ce protocole pour sa nouvelle bibliothèque numérique, la bibliothèque Cujas transforme alors le rôle de l'usager. Il n'est plus passif face à l'information transmise, mais devient un acteur de celle-ci en lui permettant à travers les outils misent à sa disposition de s'en servir pour contribuer à la valorisation des documents numérisés. Aujourd'hui, \enquote{les outils du web 2.0 ouvrent ainsi une brèche dans l'implication des utilisateurs dans la vie de la bibliothèque numérique, en leur permettant de contribuer à la valorisation et à la diffusion des collections}. L'objectif serait alors d'offrir aux utilisateurs la possibilité de contribuer à des tâches de traitement des documents, en leur permettant de faire de la transcription, ou de corriger des textes océrisés ou d'annoter les documents.

Pour répondre à ce besoin d'ouverture et de mutualisation, l'ingénierie documentaire recommande de se tourner vers les technologies du web sémantique et du protocole IIIF. L'utilisation du protocole IIIF permettrait d'alléger considérablement la charge des serveurs de la bibliothèque tout en permettant d'afficher, dans une même visionneuse experte, un volume numérise par la bibliothèque Cujas à côté d'un volume numérisé par la BNF. En utilisant le protocole IIIF, il permet de \enquote{manipuler les documents de manière décentralités}, ainsi grâce aux manifestes l'usager n'est plus \enquote{un récepteur passif}. En effet, grâce à cette outil, il est possible \enquote{d'annoter les images, comparer des feuilles à distances, et concevoir des parcours de recherche}. Le passage massif au protocole IIIF et à l'utilisation de la visionneuse Mirador ou Universal Viewer, s'observe chez la majorité des bibliothèques numériques créées aujourd'hui par les institutions patrimoniales. 

La transition vers le Web de données permettrait à la bibliothèque Cujas de s'inscrire dans un noeud dynamique au sein d'une communauté de bibliothèques entraînant une meilleure visibilité du contenu numérisé par l'établissement sur le Web, favorisant ainsi les opérations de recherche des usagers. La transition bibliographique est la traduction nationale de cette ambition, entraînant une réflexion nationale autour des codes de catalogage, comme RDA-FR, et à son éventuelle adoption.
\section{Standardisation et transition bibliographique}
La transition bibliographique a pour objectif de mieux faire remonter les données des bibliothèques dans les moteurs de recherches, par des moyens différents. La philosophie des notices bibliographies ont changé comme l'explique Marianne Clatin. \enquote{On ne catalogue plus pour que l'usager vienne sur notre catalogue; on transforme nos notices en données pour qu'elles se trouvent directement sur le chemin de l'usager, c'est-à-dire dans les moteurs de recherche généralistes (Google, etc.) grâce au référencement sémantique(SEO)}.

L'enjeu actuel pour les institutions patrimoniales, dont la bibliothèque Cujas, est d'avoir un meilleur référencement, c'est-à-dire une position parmi les premières pages retournées par Google à la suite d'une recherche. Aujourd'hui, le web a pour effet immédiat de créer une économie de l'abondance de l'information. La bibliothèque n'est donc plus un passage obligé pour accéder aux documents. Toute démarche orientée utilisateur doit donc prendre en compte comme paramètre le besoin de visibilité, permettant de capter l'attention des internautes.

\subsubsection{L'application d'IFLA-LRM aux objets complexes}
Le passage de la logique de \enquote{notices isolées} à celle de \enquote{données} interconnectées pour s'adapter au parcours de l'usager, lui permettant de faire des recherches par entités (personnes, lieux, oeuvres). Cette déconstruction de la \enquote{notice plate} au profit du modèle OEMI (Oeuvre, Expression, Manifestation et Item), décrivant et structurant les ressources bibliographiques lors de cette transition bibliographique. 

Lorsque l'on parle d'objets complexes, on désigne les périodiques représentant un objet difficile à définir dans son ensemble pour les institutions patrimoniales. Il ne s'agit pas uniquement du logiciel utilisé par les bibliothèques pour les exposer. La bibliothèque Cujas, à l'origine, a choisi de développer sa première bibliothèque numérique sous XTF avec la refonte les experts métiers envisagent une bibliothèque qui acceptent différentes types de documents. 

L'attente interne est d'intégrer les périodiques dans sa future bibliothèque. Il faut que une structure capable de \enquote{modéliser finement l'arborescence multi-niveaux propre aux revues scientifiques (Titre de la revue > Année > Fascicule/Volume > Article)}, avec une indexation détaillée poussée jusqu'au niveau de l'article. L'application des entités du modèle conceptuel IFLA-LRM permettrait de traduire cette complexité éditoriale et de combler cette lacune présente de CujasNum. Aujourd'hui, le document numérisé n'est plus un simple objet physique, mais il possède une architecture de données manipulable. Comme l'a analysé Jean-Michel Salaün, \enquote{l'utilisation de référentiels et de données modélisées en graphe pour créer du lien entre des ressources d'origine différentes, dans différents formats, qui traitent des mêmes sujets}. Ce nouveau modèle OEMI se présente avec la promesse de transformer en données interopérables tout l'information manipulée par les institutions patrimoniales. 

Les liens entre les informations se feraient à l'aide de triplet permettant à la machine de différencier les éléments grâce à l'URI, favorisant ainsi les liens de plusieurs manières des informations portant sur un même document. Il devint possible de lier une thèse de droit, par exemple, à un exemplaire physique, à un substitut numérique, à la personne qui l'a écrite, celle qui l'a dirigée en définissant précisément les relations de chacun de ces éléments. Cela permettrait de répondre à une des attentes de la direction de la BIU Cujas qui est de vouloir intégrer des réseaux nationaux en traitant de la donnée massivement. Ces liens permettent de décrire le plus précisément possible une chose (document, contenu, personne,etc.), permettantle \enquote{passeport du document}, comme l'expliquait Olesea en signalant plus il y a d'informations mieux c'est pour l'usager. L'une des solutions est donc de constituer des notices d'autorité.

\subsubsection{L'alignement d'autorités et l'URI pivot}
La majorité des notices bibliographiques sont encodés sous Dublin Core simple, qui les appauvrissent car uniquement cinq balises sur les quinze possible sont généralement utilisés. Il est important de les enrichir, notamment celles sur les livres anciens. Selon le constat d'Olesea Dubois, la richesse des informations est nécessaire pour aider l'usager à comprendre le document qu'il voit numériser aussi bien le visuel que sa matérialité. Le document numérisé est uniquement ce que voit l'usager sur la bibliothèque numérique. En effet, \enquote{la bibliothèque numérique est la seule chose que l'utilisateur voit, c'est la vitrine du document, plus on est bavard dans les descriptions}, meilleure est la compréhension intégrale du document. 

La transition d'une logique de notices isolées vers un réseau sémantique repose sur un changement dans l'art de décrire les ressources. Dans l'écosystème du Web de données, l'autorité n'est plus une chaîne de caractère figée dans un catalogue local, mais devenu une entité dynamique identifiée par une URI. Pour modéliser cette interconnexion, Emmanuelle Bermès, avec la collaboration d'Antoine Isaac et Gautier Poupeau, mobilise la métaphore du \enquote{hub and spoke}. Selon les auteurs, les référentiels, thésaurus et listes d'autorités contrôlées agissent comme le point nodal permettant de créer un ou plusieurs points de contact entre les données. \enquote{Les référentiels ou vocabulaires contrôlés qui fournissent les réservoirs d’URI réutilisables sont donc appelés à jouer un rôle essentiel dans le Web de données. Ils fournissent en effet les points d’ancrage qui vont transformer les graphes isolés de chaque institution en un graphe global, relié par des points de contacts communs}. Ces points de contact permettent de naviguer sans contrainte entre les données, en utilisant les URI. La particularité de cette métaphore est qu'il n'est pas nécessaire d'exprimer les données suivant le même modèle pour qu'elles puissent communiquer entre elles.

Appliquée à la refonte de la bibliothèque numérique de la bibliothèque Cujas, cette conceptualisation constitue un ancrage méthodologique fondamental. Elle démontre qu'il est inutile de chercher à unifier l'intégralité des structures de leurs bases de données locales ou de développer des formats propriétaires complpexes. Il suffit d'aligner les notices bibliographiques ou les fiches biographiques sur les référentiels nationaux de l'Enseignement Supérieur et de la Recherche (ESR)n principalement IdRef pour les personnes et RAMEAU pour les sujets. Sur le plan technique, cet alignement résout deux obstacles majeurs de l'accès à l'information : \enquote{l'homonymie (en distinguant deux auteurs aux patronymes identiques grâce à leur identifiant IdRef unique)}ou \enquote{la synonymie (en regroupant sous une même entité conceptuelle des termes différents ou des varitations linguistiques)}. Pour faire cet alignement de masse, deux outils existent Bibliostratus  et OpenRefine, avec des objectifs différents. 

L'alignement massif est une condition \textit{sine qua none} pour lier de manière automatisée et à grandes échelles les données locales aux grands référentiels du Web (IdRef, VIAF ou Wikidata), en brisant les silos du Web, en enrichissant les catalogues des bibliothèques sans l'intervention humaine. Un autre avantage de l'alignement est qu'il améliore l'expérience utilisateur permettant de créer des recherche par rebonds, des recherches multilinguistes et d'améliorer la visibilité des bibliothèques. Ces trois conditions sont des attentes fortes de la bibliothèque Cujas à l'issue de cette refonte qu'elle pourra proposer à ses utilisateurs. Un autre enjeu pour l'établissement est de parvenir à trouver un mode d'indexation plus adapté afin de constituer des facettes plus pertinentes pour eux également.

\subsubsection{L'alignement entre l'indexation Hauville et RAMEAU}
Les exigences de refonte sémantique formulées par le personnel de Cujas touchent également à la rationalisation de l'indexation thématique, aujourd'hui sujette à un bruit documentaire important. L'ancien responsable avait préconisé une \enquote{limitation stricte stricte à cinq indexations maximum par ouvrage}, afin de concentrer la prolifération lexicale et d'améliorer la \enquote{trouvabilité} des documents pour les chercheurs. Cette volonté de rationalisation, qui invite à se baser sur le ressenti de l'usager plutôt que d'empiler des concepts techniques complexes déconnectés des recherches réelles, rejoint la problématique de l'interopérabilité. Entre la standardisation nationale portée par le référentiel RAMEAU et la finesse de l'indexation Hauville, la négociation d'un alignement sémantique en SKOS s'impose comme une solution majeure pour répondre à ses enjeux de modernisation. Grâce à l'alignement sémantique, la bibliothèque Cujas permettrait de faire évoluer le catalogue sans tout réindexer manuelle à chaque fois qu'un terme de droit change. 

L'alignement sémantique en SKOS permettrait à la bibliothèque de conserver l'indexation Hauville, comme un thésaurus autonome et spécialisé en droit, et créer un graphe d'alignement SKOS indépendant permettant de faire le pont entre les deux. Cela réduirait alors le bruit observé lorsqu'un étudiant effectue une recherche, par exemple, dans CujasNum en cherchant à utiliser les facettes à sa disposition. 

Le choix de la bibliothèque Cujas, face à la création de sa propre indexation, résultait du constat que le référentiel RAMEAU n'est pas toujours entièrement adapté, manquant à certain moment de précision. La réfonte ne doit pas entraîner un écrasement ou l'abandon de l'indexation Hauville, mais résulte d'un travail scientifique de longue haleine. \enquote{L'harmonisation impose des choix qui peuvent entrer en conflit avec les besoins documentaires spécifiques, et pose la question de trouver comment accorder les richesses des connaissances de chaque institution concernée avec la nécessité d'avoir un minimum d'uniformité lorsqu'il s'agit de coopéreer à grande échelle}. En effet, la création d'un alignement sémantique permettrait de faire un travail de mapping entre l'indexation Hauville et le référentiel RAMEAU. En effet, comme l'a analysé Maëlys Gioan dans son étude, en expliquant qu'aligner, \enquote{c'est relier explicitement des concepts équivalents ou proches, structurer leurs correspondancess à l'aide des proprités (en SKOS, skos:axactMatch, closeMatch, broadMatch ou related)}. Cette pluralité des langages n'est pas un problèmes informatique mais une réalité humaine et scientifique, répondant à des valeurs ou des missions différentes.

L'indexation s'impose comme une brique incontournable de la constitution d'une bibliothèque numérique. Elle a besoin d'avoir fait l'objet d'un travail scientifique intellectuel conséquent car à partir de ce travai on crée les facettes. Ces facettes ont un \enquote{rôle cognitif}. Comme l'a étudié Elina Leblanc, elle explique que \enquote{les facettes reflètent la manière dont la bibliothèque numérique a pensé ses collections et aident ainsi l'utilisateur à apprendre empiriquement l'organisation des collections et à naviguer parmi les contenus}. Une indexation trop complexe ou mal hiérarchisée, comme le \enquote{problème des tirets}, perturbe cet apprentissage sémantique. Une bonne indexation est nécessaire, afin d'éviter la surabondance informatique. Rationaliser le nombre d'entrées à cinq pour l'indexation est nécessaire pour éliminer le \enquote{bruit}, réduisant la surchange informationnelle auprès des usagers.
\bigskip

En définitive, la normalisation sémantique et l'adoption de protocoles standards permet à la bibliothèque Cujas de s'assurer que ses collections patrimoniales ne sont plus de simples fichiers isolés, mais des entités interconnectées, circulant librement sur le Web de données. Pour autant, cette rigueur informatique et sémantique ne saurait constituer une fin en soi. Si le Web sémantique est conçu pour être lu et interprété par des machines, la bibliothèque numérique, elle, est crée pour des êtres humains. Les graphes de triplets (RDF) et les serveurs IIIF ne prenneent leur véritable valeur sicnetifique que lorsque cela se traduit en bénéfices concrets pour l'usager final. Toutefois, lorsque l'on conçoit une interface idéale sur le plan technique mais qui fassent abstractions des réalités cognitives, des compétences informationnelles et de la façon dont les équipes travaillent. Il est important que la future plateforme de la bibliothèque Cujas replacea l'usager au centre de la réflexion en analysant leurs attentes.